\documentclass[14pt]{extarticle}
\usepackage{fontspec}
\usepackage{polyglossia}
\usepackage{amsmath}
\usepackage{amssymb}
\usepackage{xcolor}
\usepackage{fancyhdr}
\usepackage{graphicx}
\usepackage{listings}
\usepackage[most]{tcolorbox}
% \usepackage{setspace}
\usepackage{pifont}
\usepackage{enumitem}
\usepackage{titlesec}
\usepackage[bottom]{footmisc}
\usepackage{titling}
\usepackage{minted}
\usepackage{etoolbox}
\usepackage{array}
\usepackage{longtable} % page-breakable tables (must load before polyglossia/bidi RTL setup)
\usepackage{extsizes}
\usepackage{float}
\usepackage{booktabs}
\usepackage{hyperref}
\usepackage[onehalfspacing]{setspace}
\newif\ifwhiteblack
\ifdefined\setwhiteblack
  \setwhiteblack
\fi

\ifwhiteblack
  \usepackage[a4paper, landscape, top=1.25cm, bottom=1.25cm, left=1.25cm, right=1.25cm, headheight=1.5cm, headsep=0.5cm]{geometry}
\else
  \usepackage[a4paper, top=2.7cm, bottom=2.7cm, left=2.5cm, right=2.5cm, headheight=1.5cm, headsep=0.5cm]{geometry}
\fi
\usepackage{tikz}
\usepackage{datetime2}

\usetikzlibrary{automata, positioning, arrows.meta, arrows}
\providecommand{\LTcaptype}{table}

\newfontfamily\emoji{Segoe UI Emoji}

\pagestyle{fancy}

\setmainlanguage[numerals=western]{arabic}
\setotherlanguage{english}
\newfontfamily\arabicfont[Script=Arabic]{Calibri}
\newfontfamily\arabicfonttt[Script=Arabic]{Courier New}



\ifwhiteblack
\lstset{
  language=[Sharp]C,
  numbers=left,
  stepnumber=1,
  numbersep=8pt,
  frame=single,
  basicstyle=\ttfamily\small,
  keywordstyle=\bfseries,
  stringstyle=\ttfamily,
  commentstyle=\itshape
}
\else
\lstset{
  language=[Sharp]C,
  numbers=left,
  stepnumber=1,
  numbersep=8pt,
  frame=single,
  basicstyle=\ttfamily\small,
  keywordstyle=\color{blue},
  stringstyle=\color{red},
  commentstyle=\color{green!50!black}
}
\fi

\newif\ifdetailed
\ifdefined\setdetailed
  \setdetailed
\fi

\newif\ifwithsols
\ifdefined\setwithsols
  \setwithsols
\fi

\ifwhiteblack
  \AtBeginDocument{
    \LARGE
    \usemintedstyle{bw}
  }
\fi

% Default minted look: compact code that wraps long lines.
% Not applied in black-and-white landscape mode, which already uses \LARGE.
\ifwhiteblack\else
  \setminted{fontsize=\small, breaklines=true}
\fi

% unified tcolorboxes for articles
\ifwhiteblack
\tcbset{colback=white, colframe=black, fonttitle=\bfseries, boxrule=1pt}
\newtcolorbox{boxDef}[1][]{title={{\emoji📘} تعريف\ifx\\#1\\\else ~#1\fi :}}
\newtcolorbox{boxExercise}[1][]{title={{\emoji🧩} تمرين\ifx\\#1\\\else ~#1\fi :}}
\newtcolorbox{boxExample}[1][]{title={{\emoji📝} مثال\ifx\\#1\\\else ~#1\fi :}}
\newtcolorbox{boxNote}[1][]{title={{\emoji✨} ملاحظة\ifx\\#1\\\else ~#1\fi :}}
\newtcolorbox{boxAttention}[1][]{title={{\emoji🔔} تنبيه\ifx\\#1\\\else ~#1\fi :}}
\newtcolorbox{boxWarning}[1][]{title={{\emoji⚡} ملاحظة هامة\ifx\\#1\\\else ~#1\fi :}}
\newtcolorbox{boxSolution}[1][]{title={{\emoji✅} حل\ifx\\#1\\\else ~#1\fi :}}
\newtcolorbox{boxSymbol}[1][]{title={{\emoji🔣} رمز\ifx\\#1\\\else ~#1\fi :}}
\newtcolorbox{boxHint}[1][]{title={{\emoji💡} تلميح\ifx\\#1\\\else ~#1\fi :}}
\else
\tcbset{colback=white, colframe=black, fonttitle=\bfseries, boxrule=0.8pt}
\newtcolorbox{boxDef}[1][]{colback=blue!5!white,colframe=blue!75!black,
  title={{\emoji📘} تعريف\ifx\\#1\\\else ~#1\fi :}}
\newtcolorbox{boxExercise}[1][]{colback=cyan!5!white,colframe=cyan!70!black,
  title={{\emoji🧩} تمرين\ifx\\#1\\\else ~#1\fi :}}
\newtcolorbox{boxExample}[1][]{colback=yellow!5!white,colframe=orange!90!black,
  title={{\emoji📝} مثال\ifx\\#1\\\else ~#1\fi :}}
\newtcolorbox{boxNote}[1][]{colback=gray!10!white,colframe=black,
  title={{\emoji✨} ملاحظة\ifx\\#1\\\else ~#1\fi :}}
\newtcolorbox{boxAttention}[1][]{colback=magenta!10!white,colframe=magenta!80!black,
  title={{\emoji🔔} تنبيه\ifx\\#1\\\else ~#1\fi :}}
\newtcolorbox{boxWarning}[1][]{colback=red!5!white,colframe=red!75!black,
  title={{\emoji⚡} ملاحظة هامة\ifx\\#1\\\else ~#1\fi :}}
\newtcolorbox{boxSolution}[1][]{colback=green!5!white,colframe=green!60!black,
  title={{\emoji✅} حل\ifx\\#1\\\else ~#1\fi :}}
\newtcolorbox{boxSymbol}[1][]{colback=purple!5!white,colframe=purple!70!black,
  title={{\emoji🔣} رمز\ifx\\#1\\\else ~#1\fi :}}
\newtcolorbox{boxHint}[1][]{colback=teal!5!white,colframe=teal!60!black,
  title={{\emoji💡} تلميح\ifx\\#1\\\else ~#1\fi :}}
\fi


\ifwhiteblack
\tcbset{simplecode/.style={ colback=white, colframe=black, boxrule=1pt, arc=2pt, left=4pt,right=4pt,top=4pt,bottom=4pt}}
\else
\tcbset{simplecode/.style={ colback=gray!5, colframe=black!50, boxrule=0.4pt, arc=2pt, left=4pt,right=4pt,top=4pt,bottom=4pt}}
\fi
% Compact code boxes: tighter line spacing (the preamble uses onehalfspacing) and small vertical gaps.
% Must come AFTER the simplecode definition above.
\tcbset{simplecode/.append style={fontupper=\setstretch{0.95}, top=4pt, bottom=4pt, before skip=6pt, after skip=6pt}}
\newenvironment{boxCode}{\begin{tcolorbox}[simplecode]}{\end{tcolorbox}}

% Code cell used inside tables (keeps code left-to-right inside RTL tables).
% Use inside a fixed-width p{...} column: \begin{pycell}\begin{minted}...\end{minted}\end{pycell}
\newtcolorbox{pycell}{frame empty,colback=white,boxsep=0pt,left=2pt,right=2pt,top=2pt,bottom=2pt,width=\linewidth}

\newcolumntype{C}[1]{>{\centering\arraybackslash}p{#1}}

\makeatletter
\DTMnewdatestyle{dotted}{%
  \renewcommand{\DTMdisplaydate}[4]{\two@digits{##3}.\two@digits{##2}.##1}%
  \renewcommand{\DTMDisplaydate}{\DTMdisplaydate}%
}
\makeatother
\DTMsetdatestyle{dotted}

% redefine spaces after titles
\makeatletter
\renewcommand{\@maketitle}{%
  \begin{center}
    {\huge \bfseries \@title \par}%
    \vskip 0.2em % space between title and author
    {\large \@author \par}%
    % \vskip 0.2em % space between author and date
    % {\normalsize \@date \par}%
    {\small \hfill \begin{english}\DTMtoday\end{english} \par}%
  \end{center}
}
\makeatother

\fancyhf{} % clear default
\fancypagestyle{plain}{
  \fancyhf{}
  \fancyhead[L]{مدرسة التسامح الشاملة}
  % \fancyhead[L]{\includegraphics[height=1cm]{../../../images/logoTasamoh.png}}
  \fancyhead[R]{الأستاذ محمود اغبارية}
  \fancyfoot[C]{\thepage}
}

\fancyhead[L]{مدرسة التسامح الشاملة}
\fancyhead[R]{الأستاذ محمود اغبارية}
\fancyfoot[C]{\thepage}

\setcounter{tocdepth}{3} % only section subsection and subsubsection in TOC

\newcommand{\importpdfpage}[2]{%
    \thispagestyle{empty}
    \newgeometry{left=0cm,right=0cm,top=0cm,bottom=0cm}
    \includegraphics[page=#2,width=0.95\paperwidth,height=0.95\paperheight,keepaspectratio]{#1}
    \restoregeometry
}

\newcommand{\insertFullImg}[1]{%
    \noindent
    \makebox[\textwidth][c]{
        \includegraphics[width=0.9\paperwidth,keepaspectratio]{#1}%
    }%
}%


\ifwithsols
\title{حل امتحان مفمار في علم الحاسوب للصف السابع/الثامن - بايثون $2023$}
\else
\title{امتحان مفمار في علم الحاسوب للصف السابع/الثامن - بايثون $2023$}
\fi

\begin{document}

\maketitle
\thispagestyle{fancy}

\ifwithsols
\else
\begin{boxCode}
\begin{itemize}[nosep]
    \item \textbf{الفصل الأول (إجباري - $60$ علامة):} أجب عن \textbf{جميع} الأسئلة $1$–$5$.
    \item \textbf{الفصل الثاني (اختياري - $40$ علامة):} أجب عن \textbf{سؤالين} من بين الأسئلة $6$–$9$.
    \item \textbf{مدة الامتحان:} $120$ دقيقة.
    \item \textbf{مواد مساعدة مسموحة:} كل مادة مساعدة شخصية مكتوبة أو مطبوعة، باستثناء حاسبة أو حاسوب قابل للبرمجة.
    \item اقرأ أسئلة الامتحان بتمعّن وأجب عنها بانتباه في المكان المخصّص لذلك.
\end{itemize}
\end{boxCode}
\fi

\section*{الفصل الأول — إجباري ($60$ علامة)}
في هذا الفصل خمسة أسئلة. يجب الإجابة عن \textbf{جميع} الأسئلة $1$–$5$.

\begin{enumerate}[label=\textbf{السؤال \arabic*:}, wide=0pt, leftmargin=*, itemsep=2.2em]

% ---------------------------------------------------------------
\item ($12$ علامة) أمامكم جدول فيه أربعة مقاطع برنامج بلغة بايثون. لكل مقطع برنامج في الجدول، اكتبوا ما هو المُخرج الناتج في العمود الفارغ.

{\renewcommand{\arraystretch}{1.4}
\begin{longtable}{|C{1.8cm}|p{8.7cm}|C{4cm}|}
\hline
\textbf{البند} & \centering\textbf{كود بلغة بايثون} & \textbf{المُخرج الناتج} \\
\hline
$1$. \newline ($2$ علامتان) &
\begin{pycell}
\begin{english}
\begin{minted}{python}
num1 = 5
num2 = num1 + 1
num1 = num2
print("num1 = ", num1, " num2 = ", num2)
\end{minted}
\end{english}
\end{pycell}
& \ifwithsols \begin{english}\small\texttt{num1 = 6 num2 = 6}\end{english} \fi \\
\hline
$2$. \newline ($2$ علامتان) &
\begin{pycell}
\begin{english}
\begin{minted}{python}
num = 1234
num1 = num // 100
num2 = num % 100
print("num1 = ", num1, " num2 = ", num2)
\end{minted}
\end{english}
\end{pycell}
& \ifwithsols \begin{english}\small\texttt{num1 = 12 num2 = 34}\end{english} \fi \\
\hline
$3$. \newline ($4$ علامات) &
\begin{pycell}
\begin{english}
\begin{minted}{python}
title = "Python is fun!"
print(title[3])
print(title[::2])
print(title[-2])
print(title[len(title) - 1])
\end{minted}
\end{english}
\end{pycell}
& \ifwithsols \begin{english}\texttt{h}\newline\texttt{Pto sfn}\newline\texttt{n}\newline\texttt{!}\end{english} \fi \\
\hline
$4$. \newline ($4$ علامات) &
\begin{pycell}
\begin{english}
\begin{minted}{python}
text = "Programming in python"
print(text.find('r'))
print(text.find('out'))
\end{minted}
\end{english}
\end{pycell}
& \ifwithsols \begin{english}\texttt{1}\newline\texttt{-1}\end{english} \fi \\
\hline
\end{longtable}}
\filbreak

% ---------------------------------------------------------------
\item ($12$ علامة) في الجدول الذي أمامكم أربعة مقاطع برنامج جزئية والمُخرج الذي سينتج من تشغيل مقطع البرنامج الكامل. كل مقطع برنامج يحتوي على حلقة.
أكملوا التعليمات الناقصة في كل مقطع برنامج للحصول على المُخرج المعطى. استعينوا بالمثال في أعلى الجدول.
(انتبهوا: تعليمة الطباعة مع المُعطى \textenglish{end} تضيفه إلى نهاية المُخرج بدل النزول لسطر جديد؛ مثلًا \textenglish{\texttt{print("Python", end="!")}} تطبع \textenglish{\texttt{Python!}}.)

{\renewcommand{\arraystretch}{1.4}
\begin{longtable}{|C{1.8cm}|p{8.7cm}|C{4cm}|}
\hline
\textbf{البند} & \centering\textbf{كود بلغة بايثون} & \textbf{المُخرج الناتج} \\
\hline
مثال &
\begin{pycell}
\begin{english}
\begin{minted}{python}
for item in range(3):
    print(item + 1, end=" ")
\end{minted}
\end{english}
\end{pycell}
& \begin{english}\texttt{1 2 3}\end{english} \\
\hline
$1$. \newline ($1$ علامة) &
\begin{pycell}
\begin{english}
\begin{minted}{python}
for item in range(___, ___):
    print(item)
\end{minted}
\end{english}
\end{pycell}
& \begin{english}\texttt{3}\newline\texttt{4}\end{english} \\
\hline
$2$. \newline ($2$ علامتان) &
\begin{pycell}
\begin{english}
\begin{minted}{python}
for item in range(6, ____, ____):
    print(item, end=" ")
\end{minted}
\end{english}
\end{pycell}
& \begin{english}\texttt{6 4 2}\end{english} \\
\hline
$3$. \newline ($3$ علامات) &
\begin{pycell}
\begin{english}
\begin{minted}{python}
i = 12
while i < ____:
    i = i + ____
    print(i)
\end{minted}
\end{english}
\end{pycell}
& \begin{english}\texttt{13}\newline\texttt{14}\newline\texttt{15}\end{english} \\
\hline
$4$. \newline ($6$ علامات) &
\begin{pycell}
\begin{english}
\begin{minted}{python}
i = 6
while i > ____:
    if i % 2 == ___:
        print(i)
    i = i - ___
\end{minted}
\end{english}
\end{pycell}
& \begin{english}\texttt{5}\newline\texttt{3}\newline\texttt{1}\end{english} \\
\hline
\end{longtable}}

\ifwithsols
\begin{boxSolution}
\begin{enumerate}[label=\arabic*.]
\item \textenglish{\texttt{for item in range(3, 5):}}
\item \textenglish{\texttt{for item in range(6, 0, -2):}} (أو \textenglish{\texttt{range(6, 1, -2)}})
\item \textenglish{\texttt{while i < 15:}} والسطر الذي يليه \textenglish{\texttt{i = i + 1}}
\item \textenglish{\texttt{while i > 0:}}، الشرط \textenglish{\texttt{if i \% 2 == 1:}}، والتحديث \textenglish{\texttt{i = i - 1}}
\end{enumerate}
\end{boxSolution}
\fi
\filbreak

% ---------------------------------------------------------------
\item ($12$ علامة) أمامكم عمودان من مقاطع البرنامج التي تُنفَّذ فيها حلقة. يجب الملاءمة بين أزواج الحلقات المتكافئة (بحيث ينتج في مقطعي الحلقة نفس المُخرج تمامًا) وكتابتها في الجدول أدناه.

{\renewcommand{\arraystretch}{1.3}
\begin{longtable}{|C{0.9cm}|p{5.8cm}|C{0.9cm}|p{5.8cm}|}
\hline
\multicolumn{2}{|c|}{\textbf{العمود أ}} & \multicolumn{2}{c|}{\textbf{العمود ب}} \\
\hline
$1$ &
\begin{pycell}
\begin{english}
\begin{minted}{python}
i = 1
while i < 10:
    print(i)
    i = i + 2
\end{minted}
\end{english}
\end{pycell}
& $A$ &
\begin{pycell}
\begin{english}
\begin{minted}{python}
for item in range(1, 12):
    print(item)
\end{minted}
\end{english}
\end{pycell}
\\ \hline
$2$ &
\begin{pycell}
\begin{english}
\begin{minted}{python}
for item in range(-5, 6):
    print(item + 6)
\end{minted}
\end{english}
\end{pycell}
& $B$ &
\begin{pycell}
\begin{english}
\begin{minted}{python}
i = 1
while i <= 10:
    print(i * 3)
    i = i + 1
\end{minted}
\end{english}
\end{pycell}
\\ \hline
$3$ &
\begin{pycell}
\begin{english}
\begin{minted}{python}
i = 1
while i <= 30:
    if i % 3 == 0:
        print(i)
    i = i + 1
\end{minted}
\end{english}
\end{pycell}
& $C$ &
\begin{pycell}
\begin{english}
\begin{minted}{python}
i = 8
while i > 0:
    print(i * 2)
    i = i - 2
\end{minted}
\end{english}
\end{pycell}
\\ \hline
$4$ &
\begin{pycell}
\begin{english}
\begin{minted}{python}
for i in range(16, 3, -4):
    print(i)
\end{minted}
\end{english}
\end{pycell}
& $D$ &
\begin{pycell}
\begin{english}
\begin{minted}{python}
for item in range(-4, 5, 2):
    print(item + 5)
\end{minted}
\end{english}
\end{pycell}
\\ \hline
\end{longtable}}

{\renewcommand{\arraystretch}{1.3}
\begin{longtable}{|C{5cm}|C{5cm}|}
\hline
\textbf{العمود أ ($1$ / $2$ / $3$ / $4$)} & \textbf{العمود ب ($A$ / $B$ / $C$ / $D$)} \\
\hline
$1$ & \ifwithsols $\mathbf{D}$ \fi \\ \hline
$2$ & \ifwithsols $\mathbf{A}$ \fi \\ \hline
$3$ & \ifwithsols $\mathbf{B}$ \fi \\ \hline
$4$ & \ifwithsols $\mathbf{C}$ \fi \\ \hline
\end{longtable}}
\filbreak

% ---------------------------------------------------------------
\item ($12$ علامة) في دكان \textbf{"المفاجآت"}، قرّروا منح هدية لكل زبون وفقًا لمبلغ الشراء من الدكان. يمكن أن تكون الهدية واحدة من $5$ إمكانيات:
دفتر يوميات (\textenglish{Diary})، مقلمة مجهّزة (\textenglish{Pencil case})، حقيبة مجهّزة (\textenglish{Bag})، لعبة حاسوب (\textenglish{Computer game})، قسيمة شراء بقيمة $500$ شيكل (\textenglish{Gift voucher}).
الشراء بأقل من $10$ شيكل لا يؤهّل للحصول على هدية.
تُحدَّد الهدية وفقًا للسيناريو التالي (المتغير \textenglish{amount} يمثّل مبلغ الشراء):

\begin{boxCode}
\begin{english}
\begin{minted}{python}
if amount < 10:
    print("Not eligible for a gift")
elif amount < 400:
    if amount >= 10 and amount <= 100:
        print("Diary")
    elif amount < 250:
        print("Pencil case")
    else:
        print("Bag")
elif amount <= 600:
    print("Computer game")
else:
    print("Gift voucher")
\end{minted}
\end{english}
\end{boxCode}

\begin{enumerate}[label=\alph*.]
\item ($8$ علامات) قرّر يوسف تشغيل السيناريو. ساعدوا يوسف وأكملوا الجدول: لكل مبلغ شراء، ما هي الهدية التي يستحقها الزبون؟

{\renewcommand{\arraystretch}{1.4}
\begin{longtable}{|C{3cm}|C{6cm}|}
\hline
\textbf{مبلغ الشراء} & \textbf{ما هي الهدية التي يستحقها الزبون؟} \\
\hline
$8$ & \ifwithsols لا يستحق هدية \fi \\ \hline
$80$ & \ifwithsols دفتر يوميات \fi \\ \hline
$250$ & \ifwithsols حقيبة مجهّزة \fi \\ \hline
$450$ & \ifwithsols لعبة حاسوب \fi \\ \hline
\end{longtable}}
\filbreak

\item ($4$ علامات) لائموا بين كل هدية ومجال مبالغ الشراء الملائم، واملؤوا الجدول أدناه:

{\renewcommand{\arraystretch}{1.2}
\begin{longtable}{p{5cm}p{5cm}}
$1$. دفتر يوميات & $A$. أكبر من $600$ \\
$2$. حقيبة مجهّزة & $B$. $250 - 399$ \\
$3$. لعبة حاسوب & $C$. $400 - 600$ \\
$4$. قسيمة شراء & $D$. $10 - 100$ \\
\end{longtable}}

{\renewcommand{\arraystretch}{1.3}
\begin{longtable}{|C{4.5cm}|C{4.5cm}|}
\hline
\textbf{العمود أ ($1$ / $2$ / $3$ / $4$)} & \textbf{العمود ب ($A$ / $B$ / $C$ / $D$)} \\
\hline
$1$ & \ifwithsols $\mathbf{D}$ \fi \\ \hline
$2$ & \ifwithsols $\mathbf{B}$ \fi \\ \hline
$3$ & \ifwithsols $\mathbf{C}$ \fi \\ \hline
$4$ & \ifwithsols $\mathbf{A}$ \fi \\ \hline
\end{longtable}}
\filbreak
\end{enumerate}
\filbreak

% ---------------------------------------------------------------
\item ($12$ علامة) أمامكم أربعة مقاطع برنامج. الهدف من كل مقطع برنامج هو استقبال $10$ علامات، عدّ وطباعة كمية العلامات \textbf{الأعلى من $55$ ولكن أقل من $80$}.
لسلسلة العلامات التالية: $76, 84, 80, 57, 77, 90, 68, 74, 100, 55$، يُطبع $5$.
حدّدوا لكل مقطع برنامج هل ينفّذ المطلوب في السؤال. إذا نعم، اشرحوا؛ وإذا لا، اقترحوا طريقة لتصحيح الكود لكي ينفّذ المطلوب.

\begin{enumerate}[label=\textbf{مقطع \alph*:}, wide=0pt, leftmargin=*, itemsep=1.5em]

\item .
\begin{boxCode}
\begin{english}
\begin{minted}{python}
count = 0
for item in range(10):
    grade = int(input("Please enter a grade:"))
    if grade > 55 and grade < 80:
        count = count + 1
print(count)
\end{minted}
\end{english}
\end{boxCode}
\ifwithsols
\begin{boxSolution}
\textbf{نعم، ينفّذ المطلوب.} الحلقة تتكرّر $10$ مرات، وفي كل تكرار تستقبل علامة وتفحص بواسطة العامل \textenglish{and} إذا كانت أكبر من $55$ وأصغر من $80$ في آنٍ واحد، وتزيد العداد \textenglish{count} بـ $1$، وفي النهاية تطبع العداد.
\end{boxSolution}
\else
\textbf{هل ينفّذ الكود المطلوب؟} \underline{\hspace{4cm}} \quad \textbf{الشرح / التصحيح:} \underline{\hspace{6cm}}
\fi
\filbreak

\item .
\begin{boxCode}
\begin{english}
\begin{minted}{python}
count = 0
for item in range(10):
    grade = int(input("Please enter a grade:"))
    if grade > 55 or grade < 80:
        count = count + 1
print(count)
\end{minted}
\end{english}
\end{boxCode}
\ifwithsols
\begin{boxSolution}
\textbf{لا ينفّذ المطلوب.} لأنّه استعمل العامل المنطقي \textenglish{or} بدل \textenglish{and}، والشرط \textenglish{\texttt{grade > 55 or grade < 80}} يتحقّق لأيّ عدد كان. \\
\textbf{التصحيح:} استبدال \textenglish{or} بـ \textenglish{and} في سطر الشرط: \textenglish{\texttt{if grade > 55 and grade < 80:}}.
\end{boxSolution}
\else
\textbf{هل ينفّذ الكود المطلوب؟} \underline{\hspace{4cm}} \quad \textbf{الشرح / التصحيح:} \underline{\hspace{6cm}}
\fi
\filbreak

\item .
\begin{boxCode}
\begin{english}
\begin{minted}{python}
count = 0
index = 0
while index < 10:
    grade = int(input("Please enter a grade:"))
    if grade > 55 and grade < 80:
        count = count + 1
print(count)
\end{minted}
\end{english}
\end{boxCode}
\ifwithsols
\begin{boxSolution}
\textbf{لا ينفّذ المطلوب.} لأنّ المتغير \textenglish{index} لا يتمّ تحديثه داخل حلقة \textenglish{while}، فيبقى $0$ وتدخل الحلقة في تكرار لا نهائي. \\
\textbf{التصحيح:} إضافة السطر \textenglish{\texttt{index = index + 1}} في نهاية جسم حلقة \textenglish{while}.
\end{boxSolution}
\else
\textbf{هل ينفّذ الكود المطلوب؟} \underline{\hspace{4cm}} \quad \textbf{الشرح / التصحيح:} \underline{\hspace{6cm}}
\fi
\filbreak

\item .
\begin{boxCode}
\begin{english}
\begin{minted}{python}
count = 0
index = 10
while index > 0:
    grade = int(input("Please enter a grade:"))
    if grade > 55 and grade < 80:
        print(count + 1)
    index = index - 1
print(count)
\end{minted}
\end{english}
\end{boxCode}
\ifwithsols
\begin{boxSolution}
\textbf{لا ينفّذ المطلوب.} لأنّه عند تحقق الشرط يطبع \textenglish{\texttt{count + 1}} دون أن يحدّث قيمة المتغير \textenglish{count} نفسه، فيبقى \textenglish{count} مساويًا لـ $0$ ويطبع $0$ في النهاية. \\
\textbf{التصحيح:} استبدال \textenglish{\texttt{print(count + 1)}} بـ \textenglish{\texttt{count = count + 1}}.
\end{boxSolution}
\else
\textbf{هل ينفّذ الكود المطلوب؟} \underline{\hspace{4cm}} \quad \textbf{الشرح / التصحيح:} \underline{\hspace{6cm}}
\fi
\filbreak

\end{enumerate}
\filbreak

\end{enumerate}

\clearpage
\section*{الفصل الثاني — اختياري ($40$ علامة)}
في هذا الفصل $4$ أسئلة. يجب الإجابة عن \textbf{سؤالين} من بين الأسئلة $6$–$9$ (لكل سؤال — $20$ علامة).

\begin{enumerate}[resume, label=\textbf{السؤال \arabic*:}, wide=0pt, leftmargin=*, itemsep=2.2em]

% ---------------------------------------------------------------
\item ($20$ علامة) \textbf{العدد "الخاص"} هو عدد صحيح تكون فيه منزلة الآحاد \textbf{أكبر من نصف مجموع منازل العدد}.
\begin{itemize}
\item \textbf{مثال $1$:} $129$ هو عدد "خاص" لأنّ منزلة الآحاد ($9$) أكبر من نصف مجموع منازل العدد: $1 + 2 + 9 = 12$، ونصفه $12 / 2 = 6$.
\item \textbf{مثال $2$:} $137$ هو عدد "خاص" لأنّ منزلة الآحاد ($7$) أكبر من نصف مجموع منازل العدد: $1 + 3 + 7 = 11$، ونصفه $11 / 2 = 5.5$.
\item \textbf{مثال $3$:} $31$ ليس عددًا "خاصًا" لأنّ منزلة الآحاد ($1$) ليست أكبر من نصف مجموع منازل العدد: $3 + 1 = 4$، ونصفه $4 / 2 = 2$.
\end{itemize}

\begin{enumerate}[label=\alph*.]
\item ($4$ علامات) حسب التعريف أعلاه، هل $867$ هو عدد "خاص"؟ إذا نعم، اعرضوا تعبيرًا رياضيًا (كما في المثال) يُظهر ذلك؛ وإلّا اشرحوا.
\ifwithsols
\begin{boxSolution}
\textbf{لا}، العدد $867$ ليس عددًا "خاصًا" لأنّ منزلة الآحاد هي $7$، ومجموع منازل العدد هو $8 + 6 + 7 = 21$، ونصف المجموع هو $21 / 2 = 10.5$، والمنزلة $7$ ليست أكبر من $10.5$.
\end{boxSolution}
\else
\\ \underline{\hspace{14cm}}
\fi
\filbreak

\item ($16$ علامة) أمامكم مقطع برنامج يستقبل عددًا من المستخدم ويفحص هل العدد "خاص". أكملوا الفراغات $(1)$–$(6)$ في المقطع:

\begin{boxCode}
\begin{english}
\begin{minted}{python}
def special_number():
    num = int(input("Please enter a number:"))
    total = (1)_________
    digit = num % (2)_________
    while (3)_________:
        total = total + (4)___________
        num = num // (5)_________
    if (6)________________:
        print("Number is special")
    else:
        print("Number is not special")

special_number()
\end{minted}
\end{english}
\end{boxCode}

\ifwithsols
\begin{boxSolution}
\begin{enumerate}[label=(\arabic*)]
\item $0$
\item $10$
\item \textenglish{\texttt{num > 0}}
\item \textenglish{\texttt{num \% 10}}
\item $10$
\item \textenglish{\texttt{digit > total / 2}}
\end{enumerate}
\end{boxSolution}
\fi
\filbreak
\end{enumerate}
\filbreak

% ---------------------------------------------------------------
\item ($20$ علامة)
\begin{enumerate}[label=\alph*.]
\item ($8$ علامات) أمامكم رقعة رسم فيها طول وعرض كل مربّع (خلية) هو $40$ خطوة. وفقًا للكود التالي، ارسموا على الرقعة ما سيرسمه السلحفاة (علماً أنّه يبدأ من النقطة $(0,0)$ متّجهًا لليمين):

\begin{boxCode}
\begin{english}
\begin{minted}{python}
import turtle
wn = turtle.Screen()
player = turtle.Turtle()
player.shape("turtle")
player.goto(0, 0)
player.left(180)
player.forward(80)
player.left(90)
player.forward(160)
player.left(90)
player.forward(80)
player.left(90)
player.forward(80)
player.left(90)
player.forward(80)
turtle.mainloop()
\end{minted}
\end{english}
\end{boxCode}

\begin{center}
\begin{english}
\begin{tikzpicture}[scale=0.65]
\draw[step=1cm, gray, dashed] (-4,-4) grid (4,4);
\filldraw (0,0) circle (2pt) node[above right] {\small\textbf{(0,0)}};
\ifwithsols
\draw[ultra thick, blue] (0,0) -- (-2,0) -- (-2,-4) -- (0,-4) -- (0,-2) -- (-2,-2);
\fi
\end{tikzpicture}
\end{english}
\end{center}

\ifwithsols
\begin{boxSolution}
يتّجه السلحفاة يسارًا ($180^\circ$)، يتقدّم مربّعين ($80$ خطوة) إلى $(-2,0)$، ينزل $4$ مربّعات ($160$ خطوة) إلى $(-2,-4)$، يتّجه يمينًا مربّعين إلى $(0,-4)$، يصعد مربّعين إلى $(0,-2)$، ثم يتّجه يسارًا مربّعين إلى $(-2,-2)$، فيرسم الرمز المشابه للرقم $6$ المبيّن باللون الأزرق على الشبكة.
\end{boxSolution}
\fi
\filbreak

\item ($12$ علامة) أمامكم ثلاثة مقاطع برنامج تستعمل المكتبة \textenglish{turtle}. كل مقطع برنامج يرسم شكلًا مختلفًا. لائموا بين كل كود والشكل الذي يرسمه في الجدول أدناه (افترضوا وجود أسطر التهيئة المشتركة في بداية كل كود):

{\renewcommand{\arraystretch}{1.3}
\begin{longtable}{|C{0.9cm}|p{7.2cm}|C{0.9cm}|C{5.2cm}|}
\hline
\multicolumn{2}{|c|}{\textbf{العمود أ (الكود)}} & \multicolumn{2}{c|}{\textbf{العمود ب (الشكل)}} \\
\hline
$1$ &
\begin{pycell}
\begin{english}
\begin{minted}{python}
for index in range(5):
    player.forward(60)
    player.left(360 / 5)
turtle.mainloop()
\end{minted}
\end{english}
\end{pycell}
& $A$ &
\begin{english}
\vspace{0.1cm}
\begin{tikzpicture}[scale=0.5]
\draw[thick] (0,1.2) circle (1.2cm);
\draw[thick] (0,-1.2) circle (1.2cm);
\filldraw (0,0) circle (3pt);
\end{tikzpicture}
\end{english}
\\ \hline
$2$ &
\begin{pycell}
\begin{english}
\begin{minted}{python}
for item in range(180):
    player.forward(5)
    player.left(2)
for item in range(180):
    player.forward(5)
    player.right(2)
turtle.mainloop()
\end{minted}
\end{english}
\end{pycell}
& $B$ &
\begin{english}
\vspace{0.1cm}
\begin{tikzpicture}[scale=0.55]
\draw[thick] (0,0) -- (1.5,0) -- (1.96,1.43) -- (0.75,2.31) -- (-0.46,1.43) -- cycle;
\filldraw (0,0) circle (3pt);
\end{tikzpicture}
\end{english}
\\ \hline
$3$ &
\begin{pycell}
\begin{english}
\begin{minted}{python}
for index in range(3):
    player.forward(100)
    player.left(360 / 3)
for index in range(3):
    player.forward(100)
    player.right(360 / 3)
turtle.mainloop()
\end{minted}
\end{english}
\end{pycell}
& $C$ &
\begin{english}
\vspace{0.1cm}
\begin{tikzpicture}[scale=0.5]
\draw[thick] (0,0) -- (2.4,0) -- (1.2,2.1) -- cycle;
\draw[thick] (0,0) -- (2.4,0) -- (1.2,-2.1) -- cycle;
\filldraw (0,0) circle (3pt);
\end{tikzpicture}
\end{english}
\\ \hline
\end{longtable}}

{\renewcommand{\arraystretch}{1.3}
\begin{longtable}{|C{4.5cm}|C{4.5cm}|}
\hline
\textbf{العمود أ ($1$ / $2$ / $3$)} & \textbf{العمود ب ($A$ / $B$ / $C$)} \\
\hline
$1$ & \ifwithsols $\mathbf{B}$ (مخمّس منتظم) \fi \\ \hline
$2$ & \ifwithsols $\mathbf{A}$ (دائرتان متماستان) \fi \\ \hline
$3$ & \ifwithsols $\mathbf{C}$ (مثلثان متطابقان) \fi \\ \hline
\end{longtable}}
\filbreak
\end{enumerate}
\filbreak

% ---------------------------------------------------------------
\item ($20$ علامة) \textbf{نص "النجوم"} هو نص يُبنى من نص موجود بحيث يُستبدل كل حرف في مكان (مؤشر) \textbf{فردي} في النص الأصلي برمز النجمة \textenglish{\texttt{'*'}}.
\begin{itemize}
\item \textbf{مثال:} للنص التالي \textenglish{\texttt{"Paybtchdoen"}}، ينتج نص النجوم التالي: \textenglish{\texttt{"P*y*t*h*o*n"}}.
\end{itemize}
مقطع البرنامج الذي أمامكم يستقبل نصًا \textenglish{st} من المستخدم، ويعرض نص "النجوم". أكملوا الفراغات $(1)$–$(7)$ في المقطع:

\begin{boxCode}
\begin{english}
\begin{minted}{python}
def create_stars_str():
    st = (1)____________("Please enter a string:")
    stars_str = (2)_______________________
    for index in range((3)___________):
        if (4)_________________:
            stars_str = stars_str + (5)____________
        else:
            stars_str = stars_str + (6)____________
    print((7)_______________)

create_stars_str()
\end{minted}
\end{english}
\end{boxCode}

\ifwithsols
\begin{boxSolution}
\begin{enumerate}[label=(\arabic*)]
\item \textenglish{\texttt{input}}
\item \textenglish{\texttt{""}}
\item \textenglish{\texttt{len(st)}}
\item \textenglish{\texttt{index \% 2 == 1}} (أو \textenglish{\texttt{index \% 2 == 0}} مع عكس الفرعين $(5)$ و $(6)$)
\item \textenglish{\texttt{"*"}}
\item \textenglish{\texttt{st[index]}}
\item \textenglish{\texttt{stars\_str}}
\end{enumerate}
\end{boxSolution}
\fi
\filbreak

% ---------------------------------------------------------------
\item ($20$ علامة) طُلب من الطلاب كتابة برنامج يستقبل $5$ أزواج من الأعداد الصحيحة، ويطبع كمية الأزواج التي تحتوي على عددين لهما \textbf{نفس الزوجية} (أي كلا العددين زوجيان، أو كلا العددين فرديان).
\begin{itemize}
\item \textbf{مثال:}
\begin{itemize}
\item الزوج $1$: $8, 4 \leftarrow$ زوجية متساوية (كلا العددين زوجيان).
\item الزوج $2$: $66, 7 \leftarrow$ زوجية مختلفة.
\item الزوج $3$: $3, 5 \leftarrow$ زوجية متساوية (كلا العددين فرديان).
\item الزوج $4$: $8, 1 \leftarrow$ زوجية مختلفة.
\item الزوج $5$: $5, 1 \leftarrow$ زوجية متساوية (كلا العددين فرديان).
\end{itemize}
مُخرج البرنامج سيكون $3$ لأنّ هذه هي كمية أزواج الأعداد ذات الزوجية المتساوية.
\end{itemize}
أكملوا الفراغات $(1)$–$(5)$ في المقطع لكي ينفّذ المطلوب:

\begin{boxCode}
\begin{english}
\begin{minted}{python}
count = (1)_________
for item in range((2)________):
    num1 = int(input("Please enter first number: "))
    num2 = int(input("Please enter second number: "))
    if (3)____________________________:
        count = (4)___________________
print((5)______________)
\end{minted}
\end{english}
\end{boxCode}

\ifwithsols
\begin{boxSolution}
\begin{enumerate}[label=(\arabic*)]
\item $0$
\item $5$
\item \textenglish{\texttt{num1 \% 2 == num2 \% 2}} (أو \textenglish{\texttt{(num1 + num2) \% 2 == 0}})
\item \textenglish{\texttt{count + 1}}
\item \textenglish{\texttt{count}}
\end{enumerate}
\end{boxSolution}
\fi
\filbreak

\end{enumerate}

\vspace{1cm}
\begin{flushleft}
أرجو لكم وقتًا ممتعًا.

الأستاذ محمود اغبارية.
\end{flushleft}

\end{document}
